% =============================================================================
% Режим иллюстраций — переключите ОДНУ строку (true / false) и пересоберите PDF:
%   \usenormcontrolimgtrue   — img/normcontrol/ (для нормоконтроля)
%   \usenormcontrolimgfalse  — img/ (полное качество, печать)
% =============================================================================
\newif\ifusenormcontrolimg
\usenormcontrolimgtrue


\usepackage[14pt]{extsizes}
%\usepackage{cmap} % Улучшенный поиск русских слов в полученном pdf-файле
\usepackage[T2A]{fontenc}
\usepackage{cmap}
\usepackage[utf8]{inputenc}
\usepackage[english,russian]{babel} % Языки: русский, английский
\usepackage{makecell}
\usepackage{ragged2e}
\usepackage{pdfpages}
\usepackage{amssymb}
\usepackage{amsfonts}
\renewcommand{\emph}[1]{#1}    % Делает \emph обычным текстом
\renewcommand{\textit}[1]{#1}  % Делает \textit обычным текстом
\usepackage[
  backend=biber,
  style=gost-numeric,
  sorting=none,
  language=auto,
  autolang=other
]{biblatex}
\DeclareFieldFormat*{title}{#1\isdot}       % Для полей title
\DeclareFieldFormat*{booktitle}{#1\isdot}   % Для полей booktitle
\DeclareFieldFormat*{maintitle}{#1\isdot}   % Для полей maintitle
\DeclareFieldFormat*{issuetitle}{#1\isdot}  % Для полей issuetitle
\DeclareFieldFormat*{journaltitle}{#1}      % Для полей journaltitle
\DeclareFieldFormat*{citetitle}{#1}         % Для ссылок в тексте, если используете \citetitle
\DeclareFieldFormat{labelnumberwidth}{#1\adddot}
\setlength{\biblabelsep}{0.5em}
\addbibresource{report.bib}
% ГОСТ: первый автор + название / все авторы (не «Иванов, Петров, …» перед названием).
\makeatletter
\AtBeginBibliography{%
	\toggletrue{bbx:gostbibliography}%
	\def\bbx@gost@respdelim{\setunit{\respdelim}}%
}

% Печатные источники: без URL; электронные: [Электронный ресурс] — Режим доступа: URL: … (дата обращения: …)
\AtEveryBibitem{%
	\def\bbx@gost@respdelim{\setunit{\respdelim}}%
	\clearfield{doi}%
	\ifboolexpr{%
		test {\ifentrytype{online}}%
	}{%
	}{%
		\clearfield{url}%
		\clearfield{urldate}%
	}%
	\ifboolexpr{%
		test {\ifentrytype{thesis}}%
		or test {\ifentrytype{mastersthesis}}%
		or test {\ifentrytype{phdthesis}}%
	}{%
		\clearfield{type}%
	}{}%
}
\DefineBibliographyExtras{russian}{%
	\protected\def\mkbibdate#1#2#3{%
		\mkdayzeros{#1}.\mkmonthzeros{#2}.#3%
	}%
}
\renewbibmacro*{url+urldate}{%
	\iffieldundef{url}{}{%
		\printtext{ [Электронный ресурс]}%
		\addspace\textemdash\space
		Режим доступа\addcolon\space
		URL\addcolon\space
		\url{\thefield{url}}%
		\iffieldundef{urlyear}{}{%
			\space(дата обращения\addcolon\space
			\mkdayzeros{\thefield{urlday}}.\mkmonthzeros{\thefield{urlmonth}}.\thefield{urlyear})%
		}%
	}%
}
\renewbibmacro*{doi+eprint+url+note}{%
	\ifboolexpr{test {\ifentrytype{online}}}{%
		\usebibmacro{url+urldate+note}%
	}{%
		\ifcitation
			{\iftoggle{cbx:eprint}{\usebibmacro{eprint}}{}}%
		{\iftoggle{bbx:eprint}{\usebibmacro{eprint}}{}}%
		\newunit\newblock
		\usebibmacro{url+urldate+note}%
	}%
}
% После года — пробел, не тире, перед [Электронный ресурс]
\DeclareBibliographyDriver{online}{%
  \usebibmacro{bibindex}%
  \usebibmacro{begentry}%
  \usebibmacro{heading}%
  \newunit
  \usebibmacro{author/editor}%
  \setunit*{\labelnamepunct}%
  \usebibmacro{title}%
  \setunit{\addcolondelim}%
  \usebibmacro{translation}%
  \def\bbx@gost@respdelim{\setunit{\respdelim}}%
  \usebibmacro{byauthor}%
  \setunit*{\resppunct}%
  \iflistundef{organization}
    {}
    {\setrespdelim%
     \printlist{organization}%
     \setunit*{\resppunct}}%
  \usebibmacro{credits}%
  \setunit*{\resppunct}%
  \usebibmacro{byeditor}%
  \setunit*{\resppunct}%
  \usebibmacro{bytranslator+others}%
  \newunit\newblock
  \printfield{version}%
  \newunit\newblock
  \printlist[semicolondelim]{specdata}%
  \iflistundef{location}{%
  }{%
    \newunit\newblock
    \printlist{location}%
    \setunit*{\addcomma\space}}%
  \usebibmacro{date}%
  \iffieldundef{url}{%
    \newunit\newblock
    \printupdate%
    \newunit\newblock
    \printfield{systemreq}%
    \newunit\newblock
    \usebibmacro{doi+eprint+url+note}%
  }{%
    \printupdate%
    \printfield{systemreq}%
    \setunit{\addspace}%
    \usebibmacro{doi+eprint+url+note}%
  }%
  \newunit\newblock
  \usebibmacro{addendum+pubstate}%
  \setunit{\bibpagerefpunct}\newblock
  \usebibmacro{pageref}%
  \newunit\newblock
  \usebibmacro{relatedblock}%
	\usebibmacro{finentry}}
\makeatother

% Ссылки [4, 8] — целиком на одной строке (без разрыва внутри скобок)
\DeclareDelimFormat{multicitedelim}{\addcomma\nobreakspace}
\DeclareDelimFormat{compcitedelim}{\addcomma\nobreakspace}
\makeatletter
\let\mkbibbrackets@orig\mkbibbrackets
\renewcommand*{\mkbibbrackets}[1]{\mbox{\mkbibbrackets@orig{#1}}}
\makeatother

\usepackage{titlesec}
\usepackage{float}
\usepackage{placeins}
\usepackage{enumitem} % Нумерованные и Маркированные списки
\usepackage[para,online,flushleft]{threeparttable} % Таблицы
\usepackage{array}
\usepackage{longtable}
\usepackage{etoolbox}
\setlength{\LTpre}{0pt}
\setlength{\LTpost}{0pt}
\makeatletter
% Верхнее поле 20 мм на страницах продолжения longtable
\patchcmd{\LT@output}{%
	\global\vsize\@colroom
	\copy\LT@head\nobreak}{%
	\global\vsize\@colroom
	\kern-\dimexpr 2.86mm\relax
	\copy\LT@head\nobreak}{}{}
\makeatother
\usepackage[14pt]{extsizes} % Шрифт 14 по ГОСТ 7.32-2001
\usepackage{caption} % Подписи к таблицам и рисункам
    \captionsetup{labelsep=endash} % Разделитель между номером и текстом
    \captionsetup[figure]{name={Рисунок}} % Подпись "Рисунок"
%\usepackage{amsmath} % Мат. формулы
\usepackage{amsmath}
\RequirePackage{microtype}
\sloppy

% Работа с изображениями и таблицами; переопределение названий по ГОСТу
\captionsetup[table]{singlelinecheck=false, labelsep=endash}
\captionsetup[table]{justification=raggedright,singlelinecheck=off}
\captionsetup[table]{singlelinecheck=false, labelsep=endash}

% Сжатие потоков PDF (не заменяет уменьшение исходных растров)
\pdfcompresslevel=9
\pdfobjcompresslevel=3

\usepackage{geometry} % Отступы
    \geometry{left=30mm}
    \geometry{right=15mm}
    \geometry{top=20mm}
    \geometry{bottom=20mm}
    \geometry{headheight=14pt}
    \geometry{headsep=10pt}
    \geometry{footskip=10mm}
    \setlength{\parindent}{1.25cm} % Устанавливает отступ первой строки
\raggedbottom % не растягивать вертикальные пробелы (оглавление, переносы longtable)

% Номер страницы: по центру области набора (поля 30 мм слева, 15 мм справа)
\usepackage{eso-pic}
\usepackage{fancyhdr}
\pagestyle{fancy}
\fancyhf{}
\renewcommand{\headrulewidth}{0pt}
\renewcommand{\footrulewidth}{0pt}
\fancypagestyle{plain}{%
	\fancyhf{}%
	\renewcommand{\headrulewidth}{0pt}%
	\renewcommand{\footrulewidth}{0pt}%
}
\newif\ifgostpagenumberon
\gostpagenumberonfalse
\newlength{\gostpagenumx}
\newcommand{\enablegostpagenumber}{%
	% центр текста на листе: как у \chapter* в titleformat (не 0.5\paperwidth)
	\setlength{\gostpagenumx}{\dimexpr 1in+\oddsidemargin+0.5\textwidth\relax}%
	\global\gostpagenumberontrue
}
\AddToShipoutPictureFG{%
	\ifgostpagenumberon
		\AtPageLowerLeft{%
			\put(\LenToUnit{\gostpagenumx},\LenToUnit{10mm}){%
				\makebox(0,0)[c]{\normalfont\thepage}%
			}%
		}%
	\fi
}

\usepackage{titlesec} % Используется для форматирования заголовков
\titlespacing*{\section}{0pt}{*4}{*4}
\titlespacing*{\subsection}{0pt}{*4}{*4}
\titlespacing*{\subsubsection}{0pt}{*3.5}{*2}

\usepackage{setspace} % Интервал
\onehalfspacing % Полуторный интервал

\frenchspacing % Промежутки между словами и предложениями отличаться не должны
\usepackage{indentfirst} % Красная строка

\usepackage{multirow} % Объединение строк в таблице
\usepackage{listings} % Листинги
\usepackage{xcolor} % Цвета

% \usepackage{listingsutf8}
% Для листинга кода:
\lstset{%
	language=c++,   					% выбор языка для подсветки	
	basicstyle=\small\ttfamily,			% размер и начертание шрифта для подсветки кода
	numbersep=5pt,
	numbers=left,						% где поставить нумерацию строк (слева\справа)
	%numberstyle=,					    % размер шрифта для номеров строк
	stepnumber=1,						% размер шага между двумя номерами строк
	xleftmargin=17pt,
	showstringspaces=false,
	numbersep=5pt,						% как далеко отстоят номера строк от подсвечиваемого кода
	frame=none,							% рамка через mdframed (см. ниже)
	tabsize=4,							% размер табуляции по умолчанию равен 4 пробелам
	captionpos=t,						% позиция заголовка вверху [t] или внизу [b]
	breaklines=true,					
	breakatwhitespace=true,				% переносить строки только если есть пробел
	escapeinside={\#*}{*)},				% если нужно добавить комментарии в коде
	backgroundcolor=\color{white}
}

% Листинги: подпись вне рамки (до PreInit); рамка mdframed; при переносе — сомкнутые части + «Продолжение».
\usepackage[framemethod=tikz]{mdframed}
\renewcommand{\lstlistingname}{Листинг}
\newcommand{\lstsplitcontinuation}{%
	\par\vspace{4pt}%
	\noindent{\normalfont Продолжение \lstlistingname~\thelstlisting.}%
	\par\vspace{6pt}%
}
\mdfdefinestyle{lstsplitframe}{%
	linewidth=0.8pt,
	linecolor=black,
	backgroundcolor=white,
	innertopmargin=4pt,
	innerbottommargin=4pt,
	innerleftmargin=0pt,
	innerrightmargin=0pt,
	splittopskip=18pt,
	splitbottomskip=6pt,
	middlelinecolor=black,
	middlelinewidth=0.8pt,
	nobreak=false,
	secondextra=\lstsplitcontinuation,
	middleextra=\lstsplitcontinuation,
}
% Подпись листинга вне рамки, слева (в .tex перед \begin{lstlisting}, без caption=).
\newcommand{\listingcaption}[1]{%
	\refstepcounter{lstlisting}%
	\par\vspace{6pt}%
	\noindent{\lstlistingname~\thelstlisting~--~#1}%
	\par\vspace{4pt}%
}
\surroundwithmdframed[style=lstsplitframe]{lstlisting}
\makeatletter
\lst@AddToHook{PreInit}{%
	\interlinepenalty=10000\clubpenalty=10000\widowpenalty=10000\displaywidowpenalty=10000%
}
\makeatother

\lstset{
	literate=
	{а}{{\selectfont\char224}}1
	{б}{{\selectfont\char225}}1
	{в}{{\selectfont\char226}}1
	{г}{{\selectfont\char227}}1
	{д}{{\selectfont\char228}}1
	{е}{{\selectfont\char229}}1
	{ё}{{\"e}}1
	{ж}{{\selectfont\char230}}1
	{з}{{\selectfont\char231}}1
	{и}{{\selectfont\char232}}1
	{й}{{\selectfont\char233}}1
	{к}{{\selectfont\char234}}1
	{л}{{\selectfont\char235}}1
	{м}{{\selectfont\char236}}1
	{н}{{\selectfont\char237}}1
	{о}{{\selectfont\char238}}1
	{п}{{\selectfont\char239}}1
	{р}{{\selectfont\char240}}1
	{с}{{\selectfont\char241}}1
	{т}{{\selectfont\char242}}1
	{у}{{\selectfont\char243}}1
	{ф}{{\selectfont\char244}}1
	{х}{{\selectfont\char245}}1
	{ц}{{\selectfont\char246}}1
	{ч}{{\selectfont\char247}}1
	{ш}{{\selectfont\char248}}1
	{щ}{{\selectfont\char249}}1
	{ъ}{{\selectfont\char250}}1
	{ы}{{\selectfont\char251}}1
	{ь}{{\selectfont\char252}}1
	{э}{{\selectfont\char253}}1
	{ю}{{\selectfont\char254}}1
	{я}{{\selectfont\char255}}1
	{А}{{\selectfont\char192}}1
	{Б}{{\selectfont\char193}}1
	{В}{{\selectfont\char194}}1
	{Г}{{\selectfont\char195}}1
	{Д}{{\selectfont\char196}}1
	{Е}{{\selectfont\char197}}1
	{Ё}{{\"E}}1
	{Ж}{{\selectfont\char198}}1
	{З}{{\selectfont\char199}}1
	{И}{{\selectfont\char200}}1
	{Й}{{\selectfont\char201}}1
	{К}{{\selectfont\char202}}1
	{Л}{{\selectfont\char203}}1
	{М}{{\selectfont\char204}}1
	{Н}{{\selectfont\char205}}1
	{О}{{\selectfont\char206}}1
	{П}{{\selectfont\char207}}1
	{Р}{{\selectfont\char208}}1
	{С}{{\selectfont\char209}}1
	{Т}{{\selectfont\char210}}1
	{У}{{\selectfont\char211}}1
	{Ф}{{\selectfont\char212}}1
	{Х}{{\selectfont\char213}}1
	{Ц}{{\selectfont\char214}}1
	{Ч}{{\selectfont\char215}}1
	{Ш}{{\selectfont\char216}}1
	{Щ}{{\selectfont\char217}}1
	{Ъ}{{\selectfont\char218}}1
	{Ы}{{\selectfont\char219}}1
	{Ь}{{\selectfont\char220}}1
	{Э}{{\selectfont\char221}}1
	{Ю}{{\selectfont\char222}}1
	{Я}{{\selectfont\char223}}1
}

\usepackage{pgfplots}
\pgfplotsset{compat=1.18}
\usetikzlibrary{datavisualization}
\usetikzlibrary{datavisualization.formats.functions}

\usepackage{graphicx}
% Каталог иллюстраций: imgmode.tex (\usenormcontrolimgtrue / false)
\ifusenormcontrolimg
	\def\imgroot{img/normcontrol/}
	\graphicspath{{img/normcontrol/}}
\else
	\def\imgroot{img/}
	\graphicspath{{img/}}
\fi
% \IfFileExists не смотрит \graphicspath — проверяем полный путь
\newcommand{\IfReportImgExists}[3]{%
	\IfFileExists{\imgroot #1}{#2}{#3}%
}
% Всегда из img/ или img/normcontrol/ (не из report/diags/ и т.п.)
\newcommand{\reportimg}[2][]{%
	\includegraphics[#1]{\imgroot #2}%
}

\usepackage[justification=centering]{caption} % Настройка подписей float объектов
\captionsetup[lstlisting]{%
	singlelinecheck=false,
	justification=raggedright,
	position=top,
	skip=6pt,
	labelsep=endash,
}

\usepackage[unicode,pdftex]{hyperref} % Ссылки в pdf
\hypersetup{hidelinks}

\usepackage{csvsimple}
\newcommand{\code}[1]{\texttt{#1}}
\newcommand{\termrow}[2]{%
	\textbf{#1} & {\setlength{\parindent}{0pt}#2.} \\%
	\noalign{\vspace{10pt}}%
}

\usepackage{xcolor}
\usepackage{color} 

% Главы 1--4: слева с красной строки; с главы 5 --- по центру
\titlespacing*{\chapter}{0pt}{-40pt}{8pt}
% \chapter* (реферат, содержание, определения): без отрицательного отступа --- иначе ломается верхнее поле 20 мм
\titleformat{name=\chapter,numberless}[display]{%
	\normalfont\large\bfseries
	\leavevmode
	\hbox to \textwidth{%
		\hss
		\hbox to 0pt{%
			\hss
			\hspace*{\dimexpr 0.5\paperwidth - 1in - \oddsidemargin - 0.5\textwidth\relax}%
			##1\hss
		}%
		\hss
	}%
}{}{0pt}{}
\titlespacing*{name=\chapter,numberless}{0pt}{0pt}{20pt}
\titleformat{\chapter}[block]{%
	\ifnum\value{chapter}<5\relax
		\hspace*{\parindent}\LARGE\bfseries
	\else
		\filcenter\LARGE\bfseries
	\fi
}{%
	\ifnum\value{chapter}<5\relax
		\hspace*{\parindent}\thechapter
	\else
		\thechapter
	\fi
}{20pt}{}

% section/subsection/subsubsection: красная строка + выравнивание по левому краю
\newcommand{\secheadingpre}{%
	\raggedright
	\hyphenpenalty=10000\exhyphenpenalty=10000%
	\doublehyphendemerits=10000\finalhyphendemerits=10000%
}

\titleformat{\section}[block]{\hspace*{\parindent}\Large\bfseries}{\hspace*{\parindent}\thesection}{20pt}{\secheadingpre}
\titleformat{name=\section,numberless}[block]{\hspace*{\parindent}\Large\bfseries}{}{0pt}{\secheadingpre}
\titleformat{\subsection}[block]{\hspace*{\parindent}\large\bfseries}{\hspace*{\parindent}\thesubsection}{20pt}{\secheadingpre}
\titleformat{name=\subsection,numberless}[block]{\hspace*{\parindent}\large\bfseries}{}{0pt}{\secheadingpre}
\titleformat{\subsubsection}[block]{\hspace*{\parindent}\bfseries}{}{0pt}{\secheadingpre}

\makeatletter
% Оглавление: без \vskip 0pt plus 0.2pt (иначе на стр. 7--8 верхнее поле ~21 мм).
\def\@dottedtocline#1#2#3#4#5{%
	\ifnum #1>\c@tocdepth \else
		\vskip \z@
		{\leftskip #2\relax \rightskip \@tocrmarg \parfillskip -\rightskip
			\parindent #2\relax\@afterindenttrue
			\interlinepenalty\@M
			\leavevmode
			\@tempdima #3\relax
			\advance\leftskip \@tempdima \null\nobreak\hskip -\leftskip
			{#4}\nobreak
			\leaders\hbox{$\m@th
				\mkern \@dotsep mu\hbox{.}\mkern \@dotsep
				mu$}\hfill
			\nobreak
			\hb@xt@\@pnumwidth{\hfil\normalfont\normalcolor #5%
				\kern-\p@\kern\p@}%
			\par}%
	\fi}
\renewcommand*{\l@chapter}[2]{%
	\ifnum \c@tocdepth>\m@ne
		\addpenalty{-\@highpenalty}%
		\addvspace{1em}%
		\@dottedtocline{0}{0pt}{1.5em}{\bfseries #1}{\bfseries #2}%
	\fi}
\makeatother

\renewcommand*{\arraystretch}{1}